\section{Module--Declaration Interaction}
\label{sec:module-declaration}

Sections~\ref{sec:import-graph} and~\ref{sec:theorem-premise} analyzed the module-level and declaration-level graphs independently. We now overlay the two: treating each module as a region in~$G_{\mathrm{thm}}$, we ask how well these human-designated regions correspond to the logical dependency structure. The answer bears directly on the central tension of the paper: module partitioning is a human decision (process), while declaration-level dependency is determined by logical necessity (product). Module cohesion---the degree to which a module's declarations depend on one another rather than on declarations elsewhere---measures the alignment between the two.

\subsection{Module Cohesion}
\label{sec:cohesion}

In $G_{\mathrm{import}}$, each module is an atomic vertex. In $G_{\mathrm{thm}}$, each module becomes a \emph{region}---a set of declaration-vertices that share a common source file---and each namespace becomes a larger region enclosing its constituent modules. The nesting is direct: the module-name tree~$T$ of \S\ref{sec:namespace-tree} partitions the vertices of $G_{\mathrm{thm}}$ into a two-level hierarchy (namespace $\supset$ module $\supset$ declaration), and the edges of $G_{\mathrm{thm}}$ can be classified by which boundaries they cross.

Consider, for illustration, the namespace \module{Data.Nat}, which contains (among others) two modules: \module{Data.Nat.Basic} and \module{Data.Nat.Defs}. The first module contains declarations such as \decl{add\_comm} and \decl{add\_left\_comm}; the second contains \decl{succ\_eq\_add\_one} and \decl{zero\_add}. In $G_{\mathrm{thm}}$, edges between these declarations fall into three categories: \emph{intra-module} edges (e.g., \decl{add\_left\_comm} citing \decl{add\_comm} within \module{Basic}), \emph{cross-module intra-namespace} edges (e.g., \decl{succ\_eq\_add\_one} citing \decl{add\_comm} across the two modules but within \module{Data.Nat}), and \emph{cross-namespace} edges (e.g., all four declarations citing \decl{Eq.refl}, defined in a different namespace entirely). Figure~\ref{fig:module-region} shows this structure alongside the corresponding fragment of~$T$.

\begin{figure}[H]
\centering
\begin{subfigure}[t]{0.36\textwidth}
\centering
\vspace{0pt}
% ── Directory structure (same style as Figure 1(a)) ──
\begin{forest}
  for tree={
    font=\scriptsize\ttfamily,
    text=blue!60!black,
    grow'=0,
    child anchor=west,
    parent anchor=south,
    anchor=west,
    calign=first,
    edge path={
      \noexpand\path [draw, black, \forestoption{edge}]
        (!u.south west) +(7.5pt,0) |- (.child anchor)\forestoption{edge label};
    },
    before typesetting nodes={
      if n=1 {insert before={[,phantom]}} {}
    },
    fit=band,
    before computing xy={l=12pt},
  }
[Data/
  [Nat/
    [Basic.lean
      [\textcolor{red!70!black}{add\_comm}]
      [\textcolor{red!70!black}{add\_left\_comm}]
    ]
    [Defs.lean
      [\textcolor{red!70!black}{succ\_eq\_add\_one}]
      [\textcolor{red!70!black}{zero\_add}]
    ]
  ]
]
\end{forest}
\caption{File-system hierarchy~$T$ (\textcolor{blue!60!black}{blue}), extended to the declaration level.}
\label{fig:tree-view}
\end{subfigure}%
\hfill
% ── Panel (b): Source code (top right, next to tree) ──
\begin{subfigure}[t]{0.60\textwidth}
\centering
\vspace{0pt}
\raggedright\scriptsize
\textsf{\bfseries\color{blue!60!black}Basic.lean}\\[2pt]
\ttfamily
\textbf{import} \textcolor{blue!60!black}{Init.Prelude}\\[3pt]
\textbf{theorem} \textcolor{red!70!black}{add\_comm} (n m : Nat)\\
\hspace{1.5em}: n + m = m + n := \textbf{by}\\
\hspace{1.5em}induction n <;> simp [\textcolor{red!70!black}{Eq.refl}]\\[3pt]
\textbf{theorem} \textcolor{red!70!black}{add\_left\_comm} (n m k : Nat)\\
\hspace{1.5em}: n + (m + k) = m + (n + k) := \textbf{by}\\
\hspace{1.5em}rw [\textcolor{red!70!black}{add\_comm}]; simp [\textcolor{red!70!black}{Eq.refl}]\\[6pt]
\textsf{\bfseries\color{blue!60!black}Defs.lean}\\[2pt]
\ttfamily
\textbf{import} \textcolor{blue!60!black}{Data.Nat.Basic}\\
\textbf{import} \textcolor{blue!60!black}{Init.Prelude}\\[3pt]
\textbf{theorem} \textcolor{red!70!black}{succ\_eq\_add\_one} (n : Nat)\\
\hspace{1.5em}: succ n = n + 1 := \textbf{by}\\
\hspace{1.5em}rw [\textcolor{red!70!black}{add\_comm}]; exact \textcolor{red!70!black}{Eq.refl} \_\\[3pt]
\textbf{theorem} \textcolor{red!70!black}{zero\_add} (n : Nat)\\
\hspace{1.5em}: 0 + n = n := \textbf{by}\\
\hspace{1.5em}induction n <;> exact \textcolor{red!70!black}{Eq.refl} \_
\caption{Proof source code: each \textcolor{red!70!black}{red identifier} in a proof term generates a dependency edge in~(c).}
\label{fig:source-view}
\end{subfigure}

\par\vspace{4pt}

% ── Panel (c): Isometric dependency graph (full width, bottom) ──
\begin{subfigure}[b]{\textwidth}
\centering
\begin{tikzpicture}[
  every node/.style={font=\scriptsize\ttfamily, inner sep=2pt},
  dot/.style={circle, fill=red!70!black, minimum size=3.5pt, inner sep=0pt},
  dep/.style={->, >=Stealth, red!70!black, semithick},
  depext/.style={->, >=Stealth, red!70!black, semithick, dashed},
  modimp/.style={->, >=Stealth, blue!60!black, thick},
]
  % ── Card 3 (top, back): module Defs ──
  \filldraw[fill=blue!8, draw=blue!50!black, semithick, line join=round]
    (0.2,4.7) -- (10.2,4.7) -- (11.4,5.09) -- (1.4,5.09) -- cycle;
  \filldraw[fill=blue!18, draw=blue!50!black, semithick, line join=round]
    (0.2,4.7) -- (10.2,4.7) -- (10.2,4.57) -- (0.2,4.57) -- cycle;
  \node[font=\tiny\sffamily, blue!60!black, anchor=west] at (10.4,4.63)
    {\textnormal{\module{Data.Nat.Defs}}};
  \node[dot] (sea) at (3.5,4.88) {};
  \node[font=\tiny\ttfamily, red!70!black, anchor=west] at (3.6,4.88) {succ\_eq\_add\_one};
  \node[dot] (za)  at (7.5,4.88) {};
  \node[font=\tiny\ttfamily, red!70!black, anchor=west] at (7.6,4.88) {zero\_add};

  % ── Card 2 (middle): module Basic ──
  \filldraw[fill=blue!8, draw=blue!50!black, semithick, line join=round]
    (0,2.5) -- (10.0,2.5) -- (11.2,2.89) -- (1.2,2.89) -- cycle;
  \filldraw[fill=blue!18, draw=blue!50!black, semithick, line join=round]
    (0,2.5) -- (10.0,2.5) -- (10.0,2.37) -- (0,2.37) -- cycle;
  \node[font=\tiny\sffamily, blue!60!black, anchor=west] at (10.2,2.43)
    {\textnormal{\module{Data.Nat.Basic}}};
  \node[dot] (ac)  at (3.3,2.68) {};
  \node[font=\tiny\ttfamily, red!70!black, anchor=east] at (3.2,2.68) {add\_comm};
  \node[dot] (alc) at (7.3,2.68) {};
  \node[font=\tiny\ttfamily, red!70!black, anchor=west] at (7.4,2.68) {add\_left\_comm};

  % ── Card 1 (bottom, front): module Init.Prelude ──
  \filldraw[fill=blue!8, draw=blue!50!black, semithick, line join=round]
    (1.0,0.15) -- (8.0,0.15) -- (8.9,0.46) -- (1.9,0.46) -- cycle;
  \filldraw[fill=blue!18, draw=blue!50!black, semithick, line join=round]
    (1.0,0.15) -- (8.0,0.15) -- (8.0,0.03) -- (1.0,0.03) -- cycle;
  \node[font=\tiny\sffamily, blue!60!black, anchor=west] at (8.2,0.09)
    {\textnormal{\module{Init.Prelude}}};
  \node[dot] (er)  at (4.8,0.3) {};
  \node[font=\tiny\ttfamily, red!70!black, anchor=west] at (4.9,0.3) {Eq.refl};

  % ── Namespace bracket (Data.Nat spans Cards 2–3) ──
  \draw[blue!50, semithick]
    (-0.2,2.37) -- (-0.5,2.37) -- (-0.5,5.09) -- (-0.2,5.09);
  \node[font=\tiny\sffamily, blue!60!black, rotate=90, anchor=south]
    at (-0.75,3.73) {namespace \textnormal{\module{Data.Nat}}};

  % ── Module import dots & arrows (blue = $G_{\mathrm{import}}$) ──
  \node[circle, fill=blue!60!black, minimum size=5.5pt, inner sep=0pt] (md) at (9.8, 4.88) {};
  \node[circle, fill=blue!60!black, minimum size=5.5pt, inner sep=0pt] (mb) at (10.0, 2.68) {};
  \node[circle, fill=blue!60!black, minimum size=5.5pt, inner sep=0pt] (mo) at (7.2, 0.3) {};
  \draw[modimp] (md) to[out=-85,in=85] (mb);                         % Defs → Basic
  \draw[modimp] (mb) to[out=-70,in=40] (mo);                         % Basic → other ns
  \draw[modimp] (md) to[out=-50,in=55] (mo);                         % Defs → other ns

  % ── Dependency arrows (red = mathematical logic) ──
  % Intra-module: the ONLY arrow that stays on its card
  \draw[dep] (alc) -- (ac);
  % Cross-module: Card 3 → Card 2 (external)
  \draw[depext] (sea) to[out=-88,in=92] (ac);
  % Cross-namespace: all → Eq.refl (external, piercing through card boundaries)
  \draw[depext] (ac)  to[out=-88,in=105] (er);
  \draw[depext] (alc) to[out=-88,in=60] (er);
  \draw[depext] (sea) to[out=-85,in=120] (er);
  \draw[depext] (za)  to[out=-85,in=45] (er);

\end{tikzpicture}
\caption{An isometric view of $G_{\mathrm{thm}}$ layered by module. Each \textcolor{blue!60!black}{blue card} represents a module file; \textcolor{red!70!black}{red nodes} are declarations. Solid \textcolor{red!70!black}{red arrows} show intra-module edges ($E_{\mathrm{int}}$); dashed \textcolor{red!70!black}{red arrows} show external edges ($E_{\mathrm{ext}}$) that pierce module boundaries. \textcolor{blue!60!black}{Blue arrows} on the right show the corresponding module-level imports from $G_{\mathrm{import}}$.}
\label{fig:dep-view}
\end{subfigure}
\caption{Three views of the same declarations. (a)~The file-system hierarchy and module-name tree~$T$ (\textcolor{blue!60!black}{blue}), extended to the declaration level: human decisions determine which declarations reside in which file. (b)~Proof source code: each \textcolor{red!70!black}{red identifier} in a proof term is a premise that creates a dependency edge in~(c). (c)~An isometric view of~$G_{\mathrm{thm}}$ (red arrows) layered by module (\textcolor{blue!60!black}{blue cards}); of six dependency edges, only one stays within its module---the rest pierce through the organizational layers, visualizing the low cohesion measured in this section. Together, the three panels show the full path from human file organization~(a) through proof terms~(b) to the dependency graph~(c). All declaration names have the \texttt{Nat.}\ prefix removed.}
\label{fig:module-region}
\end{figure}

\noindent
Module cohesion quantifies the proportion of a module's dependency traffic that stays internal---the strength of the innermost boundary. We now formalize this notion.

For each module~$m \in \mathcal{M}$, let $\mathcal{D}_m \subseteq \mathcal{D}$ denote the set of declarations defined in~$m$. We partition the edges of~$G_{\mathrm{thm}}$ incident to~$m$ into two classes:
\begin{align*}
  E_{\mathrm{int}}(m) &= \{(d_1, d_2) \in E_{\mathrm{thm}} \mid d_1 \in \mathcal{D}_m \text{ and } d_2 \in \mathcal{D}_m\}, \\
  E_{\mathrm{ext}}(m) &= \{(d_1, d_2) \in E_{\mathrm{thm}} \mid d_1 \in \mathcal{D}_m \oplus d_2 \in \mathcal{D}_m\},
\end{align*}
where $\oplus$ denotes exclusive or: an edge is external if exactly one endpoint belongs to~$m$.

\begin{figure}[ht]
\centering
\begin{tikzpicture}[
  every node/.style={font=\small},
  dot/.style={circle, fill=red!70!black, minimum size=4pt, inner sep=0pt},
  int/.style={->, >=Stealth, red!70!black, semithick},
  ext/.style={->, >=Stealth, red!70!black, semithick, dashed},
]
  % ── Module boundary ──
  \draw[blue!50!black, thick, dashed, rounded corners=4pt]
    (-0.5,-1.2) rectangle (3.1,1.6);
  \node[font=\small\sffamily, blue!60!black, anchor=south west] at (-0.5,1.65) {module $m$};

  % ── Internal nodes ──
  \node[dot, label={[font=\scriptsize]left:$d_1$}] (d1) at (0.3,0.8) {};
  \node[dot, label={[font=\scriptsize]above:$d_2$}] (d2) at (2.3,0.8) {};
  \node[dot, label={[font=\scriptsize]below:$d_3$}] (d3) at (1.3,-0.6) {};

  % ── External nodes ──
  \node[dot, label={[font=\scriptsize]above:$d_4$}] (d4) at (5.0,1.2) {};
  \node[dot, label={[font=\scriptsize]below:$d_5$}] (d5) at (5.0,-0.6) {};
  \node[font=\scriptsize, gray!70!black, anchor=north] at (5.0,-1.0) {other modules};

  % ── Internal edges (solid) ──
  \draw[int] (d1) -- (d2);
  \draw[int] (d3) -- (d1);
  \node[font=\tiny, red!70!black, anchor=south] at (1.3,0.85) {$E_{\mathrm{int}}$};

  % ── External edges (dashed) ──
  \draw[ext] (d1) -- (d4);
  \draw[ext] (d2) -- (d5);
  \draw[ext] (d3) -- (d4);
  \draw[ext] (d3) -- (d5);
  \node[font=\tiny, red!70!black, anchor=south] at (3.9,0.85) {$E_{\mathrm{ext}}$};

\end{tikzpicture}
\caption{Internal vs.\ external edges. Internal edges (\textcolor{red!70!black}{solid red}) connect declarations within the same module; external edges (\textcolor{red!70!black}{dashed red}) cross the module boundary.}
\label{fig:cohesion-schematic}
\end{figure}

\begin{definition}[Module cohesion]
\label{def:cohesion}
The \emph{cohesion} of a module~$m$ is
\[
  \mathrm{coh}(m) = \frac{|E_{\mathrm{int}}(m)|}{|E_{\mathrm{int}}(m)| + |E_{\mathrm{ext}}(m)|},
\]
with the convention that $\mathrm{coh}(m) = 0$ when $|E_{\mathrm{int}}(m)| + |E_{\mathrm{ext}}(m)| = 0$.
\end{definition}

\begin{example}
For the module~$m$ in Figure~\ref{fig:cohesion-schematic}, there are $|E_{\mathrm{int}}(m)| = 2$ internal edges (solid arrows) and $|E_{\mathrm{ext}}(m)| = 4$ external edges (dashed arrows), giving
\[
  \mathrm{coh}(m) = \frac{2}{2+4} = \frac{1}{3} \approx 0.33.
\]
One-third of the dependency traffic involving~$m$ stays internal. In \module{Mathlib}, the median module has $\mathrm{coh}(m) = 0$: \emph{no} internal traffic at all.
\end{example}

Cohesion ranges from~$0$ (no internal edges: every dependency crosses the module boundary) to~$1$ (a perfectly self-contained module with no external dependencies). The illustrative module of Figure~\ref{fig:module-region} has $\mathrm{coh} \approx 0.22$---already far higher than a typical \module{Mathlib} module, as we shall see. If module boundaries reflected the logical structure faithfully, one would expect cohesion values to cluster near~$1$: declarations within the same module would cite each other frequently and depend on external declarations sparingly.

\begin{table}[ht]
\centering
\caption{Global distribution of module cohesion across $6{,}993$ file-level modules.}
\label{tab:cohesion-global}
\begin{tabular}{lr}
\toprule
Statistic & Value \\
\midrule
Mean       & $0.107$ \\
Median     & $0.074$ \\
Std Dev    & $0.121$ \\
Max        & $1.000$ \\
Modules with $\mathrm{coh} = 0$ & $584$ / $6{,}993$ ($8.4\%$) \\
\bottomrule
\end{tabular}
\end{table}

The central finding is that module cohesion is low throughout: the mean is~$0.107$ and the median is~$0.074$, indicating that internal edges constitute only about $10\%$ of a typical module's total connectivity. While $8.4\%$ of modules have zero internal edges, even the most cohesive modules rarely approach~$1$.

These numbers have a precise structural interpretation. If module boundaries were aligned with the logical dependency structure---if ``belonging to the same module'' predicted ``depending on each other''---then cohesion values would cluster near~$1$. The observed mean of~$0.107$ indicates that the module boundary and the dependency structure are only weakly aligned: knowing that two declarations share a source file tells you little about whether one cites the other.

\begin{remark}[Levels of granularity and the namespace dependency graph]
\label{rem:ns-vs-file-cohesion}
This paper analyzes two dependency graphs: the module-level import graph~$G_{\mathrm{import}}$ (\S\ref{sec:import-graph}), whose nodes are source files, and the declaration-level premise graph~$G_{\mathrm{thm}}$ (\S\ref{sec:theorem-premise}), whose nodes are individual theorems and definitions. A third natural construction---the \emph{namespace dependency graph}~$G_{\mathrm{ns}}$, in which each namespace is a node and an edge $A \to B$ indicates that some declaration in namespace~$A$ cites a declaration in namespace~$B$---occupies an intermediate level of granularity. We do not analyze~$G_{\mathrm{ns}}$ here because the relationship between namespaces and source files is many-to-many: a single namespace may span dozens of files, and a single file may contribute declarations to multiple namespaces. This many-to-many mapping makes it difficult to disentangle the namespace-level signal from file-level artifacts without additional normalization. The construction and analysis of~$G_{\mathrm{ns}}$ is deferred to future work (\S\ref{sec:future}).

As a concrete illustration of this entanglement: if cohesion is computed using declaration-name prefixes as a proxy for modules (yielding approximately $15{,}000$ namespace-level units), the mean cohesion drops to~$0.034$ and $56\%$ of units have zero cohesion. The discrepancy arises because name-prefix grouping creates many single-declaration ``modules'' that trivially have no internal edges. The file-level cohesion reported here ($6{,}993$ modules, mean~$0.107$) is the more meaningful measure, since it corresponds to the actual organizational units that developers create and maintain.
\end{remark}

\subsection{Import Edges vs.\ Declaration-Level Dependencies}
\label{sec:import_vs_declaration}

Section~\ref{sec:import-graph} analyzed~$G_{\mathrm{import}}$, the graph of developer-written \texttt{import} statements; \S\ref{sec:theorem-premise} analyzed~$G_{\mathrm{thm}}$, the declaration-level premise graph. We now connect the two by aggregating~$G_{\mathrm{thm}}$ to the module level. Define the \emph{file-aggregated graph}~$G_{\mathrm{file}}$: its vertices are the same source files as~$G_{\mathrm{import}}$, and it contains an edge $m_1 \to m_2$ whenever some declaration in file~$m_1$ cites a declaration in file~$m_2$ (with $m_1 \neq m_2$). Comparing~$G_{\mathrm{file}}$ with~$G_{\mathrm{import}}$ reveals how much of the human-declared import structure corresponds to actual declaration-level usage.

The scale difference is striking: $G_{\mathrm{file}}$ has $215{,}211$ edges, while~$G_{\mathrm{import}}$ has only $23{,}235$---the actual module-to-module dependencies at the declaration level outnumber the human-written imports by a factor of~$9.1\times$. Table~\ref{tab:import-vs-file} decomposes this gap.

\begin{table}[ht]
\centering
\caption{Edge-level comparison of $G_{\mathrm{import}}$ (human-written imports) and $G_{\mathrm{file}}$ (declaration-aggregated dependencies).}
\label{tab:import-vs-file}
\begin{tabular}{lrr}
\toprule
Category & Edges & Percentage \\
\midrule
\multicolumn{3}{l}{\emph{$G_{\mathrm{import}}$ breakdown ($23{,}235$ edges):}} \\
\quad Active imports (in both $G_{\mathrm{import}}$ and $G_{\mathrm{file}}$) & $16{,}742$ & $72.1\%$ \\
\quad Unused imports ($G_{\mathrm{import}}$ only) & $6{,}493$ & $27.9\%$ \\
\midrule
\multicolumn{3}{l}{\emph{$G_{\mathrm{file}}$ breakdown ($215{,}211$ edges):}} \\
\quad Direct import exists & $16{,}742$ & $7.8\%$ \\
\quad Indirect, transitively reachable & $197{,}639$ & $91.8\%$ \\
\quad Indirect, not reachable & $830$ & $0.4\%$ \\
\bottomrule
\end{tabular}
\end{table}

The table tells two complementary stories. From the import side, $72.1\%$ of import edges are \emph{active}---the importing file's declarations actually cite declarations in the imported file. The remaining $27.9\%$ are \emph{unused} at the declaration level: no declaration in the importing file directly references the imported module's contents. These imports exist to provide transitive visibility, bringing into scope not the imported module itself but the modules it in turn imports.

From the declaration side, only $7.8\%$ of cross-file declaration dependencies correspond to a direct import edge. The remaining $92.2\%$ are \emph{indirect}: file~$A$'s declarations cite file~$B$'s declarations, but $A$ does not import~$B$ directly. Almost all of these ($197{,}639$ out of $198{,}469$) are transitively reachable through the import chain, confirming that the import graph, while sparse, provides sufficient transitive coverage. The $830$ edges ($0.4\%$) that are not transitively reachable correspond to cross-library references and mapping gaps at the boundary of the \module{Mathlib} source tree.

The import graph thus functions as a sparse gateway: each file declares a handful of direct imports (mean out-degree $3.1$ in~$G_{\mathrm{import}}$), but through transitivity these few edges open access to a vast dependency space (mean out-degree $30.8$ in~$G_{\mathrm{file}}$). This $10\times$ amplification quantifies the indirectness between human-level organization and declaration-level logic: the cognitive act of writing an \texttt{import} statement is a coarse, high-level gesture that implicitly enables hundreds of fine-grained logical dependencies.

\begin{remark}[Relationship to transitive redundancy]
\label{rem:unused-vs-redundant}
Section~\ref{sec:transitive-reduction} found that $17.5\%$ of import edges are \emph{transitively redundant}---removable without disconnecting the import graph. Here we find that $27.9\%$ are \emph{unused}---no declaration in the importing file directly cites the imported module. These two numbers measure different properties: redundancy is a graph-theoretic property (edge removability), while ``unused'' is a semantic property (absence of declaration-level citation). The gap between them ($27.9\% > 17.5\%$) has a precise interpretation: $78.4\%$ of unused imports are graph-theoretically \emph{essential}---they are the sole path through which some other module's declarations become transitively visible---while $21.6\%$ are both unused and transitively redundant. Conversely, some active imports are nonetheless transitively redundant: the importing file cites the imported module, but the same module was already reachable through an alternative import chain.
\end{remark}

\subsection{Cohesion by Directory}
\label{sec:cohesion-namespace}

We group the~$6{,}993$ modules by their top-level directory under \module{Mathlib/} and compute the average cohesion within each group. Table~\ref{tab:cohesion-namespace} reports the results for directories with at least~$100$ modules.

\begin{table}[ht]
\centering
\caption{Average module cohesion by top-level directory (directories with $\ge 100$ modules).}
\label{tab:cohesion-namespace}
\begin{tabular}{lrr}
\toprule
Directory & Modules & Mean cohesion \\
\midrule
\module{Combinatorics}      & $163$     & $0.188$ \\
\module{Data}               & $608$     & $0.174$ \\
\module{Logic}              & $55$      & $0.170$ \\
\module{Order}              & $293$     & $0.150$ \\
\module{GroupTheory}        & $151$     & $0.133$ \\
\module{CategoryTheory}     & $921$     & $0.115$ \\
\module{AlgebraicTopology}  & $109$     & $0.107$ \\
\module{AlgebraicGeometry}  & $105$     & $0.103$ \\
\module{Topology}           & $591$     & $0.100$ \\
\module{Algebra}            & $1{,}218$ & $0.090$ \\
\module{NumberTheory}       & $207$     & $0.079$ \\
\module{Probability}        & $111$     & $0.078$ \\
\module{LinearAlgebra}      & $325$     & $0.076$ \\
\module{Analysis}           & $702$     & $0.069$ \\
\module{RingTheory}         & $621$     & $0.068$ \\
\module{Geometry}           & $120$     & $0.066$ \\
\module{MeasureTheory}      & $284$     & $0.066$ \\
\bottomrule
\end{tabular}
\end{table}

The variation across directories spans a $2.8\times$ range, from \module{Combinatorics} ($0.188$) to \module{MeasureTheory} ($0.066$). The highest-cohesion directories---\module{Combinatorics}, \module{Data}, and \module{Order}---contain modules that develop concrete structures (graphs, lists, partial orders) with self-contained APIs: each file defines a type and proves lemmas about it, producing a comparatively high ratio of internal edges. \module{CategoryTheory}, which \S\ref{sec:thm-community} identified as the most autonomous major discipline in~$G_{\mathrm{thm}}$ ($97.4\%$ community--namespace concentration), has above-average cohesion ($0.115$), carrying its disciplinary autonomy down to the module level.

The lowest cohesion values belong to \module{MeasureTheory} ($0.066$), \module{Geometry} ($0.066$), and \module{RingTheory} ($0.068$). These directories develop theories that draw heavily on definitions from \module{Algebra}, \module{Topology}, and \module{Analysis}---their declarations are boundary-spanning by nature, reaching outward to the shared vocabulary of formalized mathematics.

\subsection{High-Cohesion and Zero-Cohesion Modules}
\label{sec:cohesion-extremes}

To identify modules whose internal structure is unusually tight or unusually sparse, we restrict attention to modules with at least~$10$ declarations (filtering out trivially small modules) and rank by cohesion.

\begin{table}[ht]
\centering
\caption{Top~10 modules by cohesion (among modules with $\ge 10$ declarations).}
\label{tab:cohesion-top}
\small
\begin{tabular}{lrrrr}
\toprule
Module & Decl.\ & Int.\ edges & Ext.\ edges & Cohesion \\
\midrule
\module{Tactic.DeriveTraversable}     & $25$  & $34$  & $0$   & $1.000$ \\
\module{Util.Notation3}               & $31$  & $40$  & $0$   & $1.000$ \\
\module{Tactic.Sat.FromLRAT}          & $34$  & $107$ & $1$   & $0.991$ \\
\module{Data.Erased}                  & $22$  & $64$  & $1$   & $0.985$ \\
\module{Tactic.Simps.Basic}           & $24$  & $38$  & $2$   & $0.950$ \\
\module{SetTheory.Lists}              & $43$  & $130$ & $11$  & $0.922$ \\
\module{Data.TwoPointing}             & $23$  & $33$  & $5$   & $0.868$ \\
\module{Combinatorics.Graph.Basic}    & $59$  & $171$ & $32$  & $0.842$ \\
\module{Data.List.DropRight}          & $39$  & $52$  & $10$  & $0.839$ \\
\module{Data.List.SplitOn}            & $12$  & $13$  & $1$   & $0.929$ \\
\bottomrule
\end{tabular}
\end{table}

The highest-cohesion modules fall into two categories. The first is \emph{tactic and metaprogramming modules}: \module{Tactic.DeriveTraversable} and \module{Util.Notation3} achieve perfect cohesion ($1.000$) because their declarations form closed internal chains with zero external edges---they are self-contained programs, not mathematical developments. The second category is \emph{concrete data structure modules}: \module{SetTheory.Lists} ($0.922$, $43$ declarations) develops a self-contained theory of hereditary lists; \module{Combinatorics.Graph.Basic} ($0.842$, $59$ declarations) defines simple graphs and proves basic properties; \module{Data.Erased} ($0.985$) implements a computationally erased wrapper type. In both categories, the file contains a type definition together with its API lemmas, producing a high ratio of internal edges.

At the opposite extreme, $584$ modules ($8.4\%$) have cohesion exactly zero. Among the zero-cohesion modules with at least~$10$ declarations, typical examples include \module{Algebra.Group.Nat.Even} ($20$ declarations, $49$ external edges), \module{Data.Set.Order} ($17$ declarations, $123$ external edges), and \module{Order.Module.Pointwise} ($16$ declarations, $176$ external edges). These modules collect lemmas about a specific predicate or construction: every lemma in \module{Algebra.Group.Nat.Even} characterizes even-number properties, but no lemma in that module cites another lemma in the same module---they all reach outward to foundational definitions elsewhere.

\begin{remark}
\label{rem:zero-cohesion}
A zero-cohesion module is not necessarily poorly organized. It is a \emph{lemma collection}: a set of declarations grouped by the concept they describe rather than by logical interdependence. Developers place ``all lemmas about even natural numbers'' in \module{Algebra.Group.Nat.Even} so that a human can find them by topic, even though these lemmas depend on definitions and lemmas scattered across many other modules. This is precisely the organizational pattern one would expect from a library designed for human navigation. Zero cohesion is not a defect; it is the structural signature of cognitive organization.
\end{remark}

\subsection{The Module as Cognitive Container}
\label{sec:cognitive-container}

We can now assemble the evidence from \S\ref{sec:import-graph}, \S\ref{sec:theorem-premise}, and the present section into a unified picture of the module's structural role. Three numbers, drawn from three levels of analysis, tell a convergent story:

\begin{table}[ht]
\centering
\caption{Cross-level comparison of module boundary effects.}
\label{tab:cross-level}
\begin{tabular}{llr}
\toprule
Section & Metric & Value \\
\midrule
\S\ref{sec:cross-namespace} & Intra-namespace edges in $G_{\mathrm{import}}^{-}$ & $62.9\%$ \\
\S\ref{sec:thm-cross-namespace} & Intra-module edges in $G_{\mathrm{thm}}$ & $9.6\%$ \\
\S\ref{sec:cohesion} & Mean module cohesion & $0.107$ \\
\bottomrule
\end{tabular}
\end{table}

At the module level, $62.9\%$ of reduced import edges stay within the same namespace (\S\ref{sec:cross-namespace})---the cognitive categories of the tree~$T$ exert a moderate organizing force on inter-module dependencies. At the declaration level, only $9.6\%$ of edges in~$G_{\mathrm{thm}}$ remain within the same module (\S\ref{sec:thm-cross-namespace})---the module boundary is nearly invisible to individual declarations. The cohesion figure of~$0.107$ confirms this from the complementary direction: within a module, the ratio of internal to total edges is small.

These three numbers describe the same phenomenon at increasing magnification. At the coarsest level, modules cluster within namespaces---the cognitive taxonomy has real organizational force. Zoom in to individual declarations, and the module boundary nearly vanishes: declarations overwhelmingly depend on declarations elsewhere. Zoom in further, asking not just about individual edges but about the \emph{internal density} of each module, and the signal drops to~$10.7\%$. The module is a container that organizes the library for human cognition---for browsing, for discovery, for maintaining a mental model of what exists where---but it is not a unit of logical structure. The declarations inside a module are only modestly more likely to depend on each other than on declarations in any other module of comparable size.

\begin{remark}
\label{rem:module-role}
The distinction between cognitive container and logical unit resolves an apparent paradox. On one hand, \module{Mathlib}'s module structure is carefully maintained: refactoring PRs regularly move files between directories, rename namespaces, and reorganize the hierarchy. On the other hand, our data shows that module boundaries have almost no relationship to logical dependency structure. Both observations are correct because they describe different functions. The module structure is maintained precisely because it serves cognitive navigation---developers need to find lemmas by topic---not because it reflects logical architecture. The maintenance effort is real but serves a human purpose, not a logical one.
\end{remark}

For \module{Mathlib} maintainers, module cohesion offers a quantitative refactoring signal: modules with cohesion near zero may be candidates for reorganization, while the low overall mean ($0.107$) suggests that a systematic audit of module boundaries could improve the alignment between file organization and logical structure.

\subsection{Implications for AI-Assisted Formalization}
\label{sec:ai-implications}

The current module structure of \module{Mathlib} is optimized for human cognition: declarations are grouped by conceptual topic so that a developer searching for ``lemmas about sorted lists'' can navigate to \module{Data.List.Sort} and find them collected in one place. This organization serves the human prover's need to browse, discover, and maintain a mental model of the library.

If AI systems become the primary producers or consumers of formal proofs, the cognitive navigation function of modules will diminish in importance. An AI premise-selection model does not browse the module-name tree; it queries a learned embedding space. An AI proof search does not rely on umbrella imports to bring relevant lemmas into scope; it retrieves premises from a global index. The organizational principles that justify the current module structure---conceptual proximity, navigable hierarchy, thematic coherence---are optimized for a cognitive agent that the AI is not.

Two divergent evolutionary paths are plausible. In one trajectory, modules degrade into pure \emph{compilation units}: their boundaries are drawn to minimize build time and maximize parallelism, with no pretense of thematic coherence. Declarations would be distributed across files by dependency depth and compilation cost rather than by mathematical topic. In the other trajectory, modules evolve into \emph{semantic interfaces}---a new kind of cognitive container that serves both human understanding and AI retrieval. Rather than grouping ``all lemmas about $X$,'' a module might group ``all declarations needed to establish interface $Y$,'' where the interface is defined by a coherent set of types and their properties. This second trajectory would likely produce modules with higher cohesion than the current~$0.107$, since interface-oriented organization aligns more closely with dependency structure.

The cohesion data of this section provides a baseline for tracking this evolution. The current mean cohesion of~$0.107$ quantifies the degree to which human cognitive organization diverges from logical structure. As AI contributions to \module{Mathlib} increase---whether through tactic automation, premise selection, or full proof generation---monitoring cohesion over time will reveal whether the module structure is drifting toward compilation efficiency, toward semantic coherence, or remaining frozen in its current cognitive-container form.